\Needspace{10\baselineskip}
\section{Source audit, normalization safeguards, and verification}
\label{audit:section}

\subsection{Inspected baseline and source attribution}

The source anchor is the ProveIt revision
\src{7bd45777a8a127e5dc69aff8887ca36a79a3059d}, including the eleven
archives then present in \src{docs/incoming} \cite{Repository}.
The archive inventory, immutable source URLs, and hashes accompany
the delivery.  The mathematical review concentrated on the relative
Hurwitz definitions, nonpositive-index reductions, one-free-order
generator, centered moment formulas, and the explicit R1, R7, and R8
questions of the incoming harmonic-parity report
\cite{IncomingParity,IncomingIndependent}.  This is a targeted analytic
review with a complete archive inventory; it is not a line-by-line
certification of every proof in the entire repository.

No newly confirmed false theorem in the inspected source material
emerged from this review.  The refinements below specify hypotheses and
normalizations that must survive integration.  In particular, the
source's warning against differentiating a fixed deleted index was
correct: Section~\ref{tr:section} supplies the extra moving-parameter
identity needed to cross that boundary.

The kernel attribution deserves explicit preservation.  The geometric
kernel reorganized in \eqref{zbg:kernel}, and the original entire
interpolation that it represents, are established in Ihara, Nakamura,
and Yamamoto \cite[Proposition~3.3 and Theorem~4.1]{INY}.  The published
reference is \emph{The Ramanujan Journal} \textbf{67} (2025), article~1,
doi:10.1007/s11139-025-01045-2; the arXiv version remains a useful
accessible precursor.  Bibliographic entries that describe this work
only as a preprint can therefore be updated.  The incoming
independent-orders report already identifies the interpolant
\cite{IncomingIndependent}.  Our contribution is the stated joint
twist completion, the all-block unit-polydisc theorem, and the
transverse formulas, rather than a new discovery of their underlying
geometric interpolation.

Likewise, the finite gap polynomials and the one-free-order generating
function retain their attribution to \cite{IncomingParity}.  Euler's
odd-weight double-zeta reduction remains a classical ingredient,
with the explicit primary proof cited in \cite{Bradley2007}.
Several incoming reports have overlapping parity and quartic-zeta
results from earlier common baselines.  Their separate proofs may be
useful, but the canonical question ledger should count the corresponding
resolved question once and retain each report's actual scope.

\subsection{An absolutely convergent boundary series can have a different value}
\label{audit:boundary}

The initial convergence domain in Theorem~\ref{tr:jointD} is essential.
An isolated exponent can cancel the leading asymptotic term of a
paired summand, producing an ordinary convergent series whose sum differs
from the continuation of the original family.  Here is a complete
elementary example, requiring no numerical regularization.

Take $m=r=0$, $a=2$, and $b=1$.  Then
\[
 \ell_{0,0}(x)=\zeta(0,x)=\frac12-x,
 \qquad h_{2,1}(s)=1.
\]
On the initial half-plane $\Re s>2$, the defining weighted difference
is
\begin{align}
 D_{0,0}(s;2,1)
 &=\sum_{n\geq0}\{f_s(n+1)-f_s(n+2)\},
 & f_s(x)&=\frac12x^{-s}-x^{1-s}.
 \label{audit:telescoping}
\end{align}
Its $N$th partial sum is exactly
\begin{equation}\label{audit:boundary-partial}
 -\frac12-\frac12(N+1)^{-s}+(N+1)^{1-s}.
\end{equation}
Consequently its entire continuation is the constant
\begin{equation}\label{audit:continued}
                         D_{0,0}(s;2,1)=-\frac12.
\end{equation}
But substituting $s=1$ in the individual paired summands first gives
the ordinary absolutely convergent series
\begin{equation}\label{audit:literal}
 \sum_{n\geq0}\frac12
       \left(\frac1{n+1}-\frac1{n+2}\right)=\frac12.
\end{equation}
The values differ by one, which is the boundary term
$(N+1)^{1-s}$ in \eqref{audit:boundary-partial}.  This discrepancy
does not contradict analytic continuation: the paired series lacks
locally normal convergence on a complex neighborhood of $s=1$.

Thus \eqref{tr:Didentity} and \eqref{tr:StieltjesD} must retain the
phrase ``the continuation from the stated initial domain.''  Literal
summation at an accidentally improved boundary point is a different
operation.  The example is also a useful negative control for software
that evaluates parameterized infinite sums after specialization.

\subsection{Hypotheses and constants to preserve during integration}

\paragraph{Moving and fixed indices.}
The polynomial identity for a word with a letter fixed at $u=-m$
contains no information about derivatives in the transverse variable
$u$.  It can be differentiated in its independently free letters.
To differentiate in $u$, first use the entire identity
\eqref{tr:genericgap}, and then specialize.  The functions
$\ell_{m,r}$ are the additional data introduced by that operation.
The finite-family bound in Proposition~\ref{tr:finitefamily} concerns
one marked index; it does not establish the same bound for several
independently marked indices.

\paragraph{Strict endpoints and Bernoulli signs.}
For a strict unit-lattice gap the finite power sum is
\[
 \sum_{z<x<y}x^m
   =\frac{B_{m+1}(y)-B_{m+1}(z+1)}{m+1}.
\]
Accordingly $B_1=-1/2$ and $B_1(1)=1/2$ occur in different positions.
At $m=0$, the missing endpoints account for the last term in
\[
 \HH(s,0,t)=\HH(s-1,t)-\HH(s,t-1)-\HH(s,t).
\]
The analogous term in \eqref{tr:commutator} must also remain after
transverse differentiation.  Its sign is fixed by the Hurwitz shift
identity, not by a choice of regularization.

\paragraph{Centered cutoff constants and the two zero-order zeta values.}
The mixed-tail proof uses the cutoff $0\leq n<N$ and the constant
term in the variable $N$.  The equality with the spectral value is
proved in \eqref{mt:cutoff-value}; it is not a general rule for every
cutoff coordinate.  Two different zero-order zeta values appear:
\[
                \zeta(0,1/2)=0,\qquad \zeta(0)=-\frac12.
\]
The first makes the continued sums of the subtracted centered even
polynomials vanish.  The second encodes the surviving rational boundary
term in \eqref{mt:boundary-constant}.  Dropping either distinction
changes the moment formula.  At $d=a+b-1$, the two oriented
$\gamma/2$ terms cancel the diagonal $-\gamma$; an isolated
$\zeta(1)$ is never introduced.

\paragraph{Domains, branches, and spanning statements.}
Endpoint parameters stay in the right half-plane.  The twists satisfy
the consecutive-sum inequalities in \eqref{zbg:tube}; removable zero
collisions do not remove the remaining nonzero-period hyperplanes.
The polylogarithmic primitives are first fixed in the right half-plane
and then continued as their complete anchored expressions, as in
\eqref{mt:L-polylog}.  The finite list in \eqref{mt:finite-span} is
a proved rational spanning list.  Its at-most-$\max(a,b)$ count is not
an assertion of arithmetic independence or minimality.  The ordinary
moment theorem also does not imply ordinary-zeta closure of its first
spectral derivative, which contains additional global logarithmic data.

\subsection{Reproducible verification and the limits of its evidence}

The analytic proofs above establish the unrestricted identities.
The companion scripts provide finite exact checks and independent
high-precision diagnostics of the most sensitive formulas.  The
recorded coverage is summarized below.

\begin{center}
\small
\begin{tabular}{@{}>{\raggedright\arraybackslash}p{1.40in}rr>{\raggedright\arraybackslash}p{3.43in}@{}}
\toprule
Family & Exact & Numerical & What is compared\\
\midrule
Geometric kernels and zero blocks
 & 590 & 0
 & Endpoint composition, deconcatenation, kernel recurrence, finite
   intervals, inverse intervals, gap coefficients, and insertion sums.\\[4pt]
Mixed-tail moments
 & 460 & 30
 & 400 boundary identities and 60 finite summation-by-parts identities;
   direct Hurwitz heads and centered tail expansions at two cutoffs.\\[4pt]
Transverse identities
 & 0 & 15
 & Finite Gamma/Hurwitz gaps, noninteger endpoint tails, Barnes values,
   negative-outer-order formulas, and the boundary discrepancy above.\\[4pt]
Normalized Mellin derivatives
 & 0 & 24
 & Subtracted logarithmic quadrature versus derivatives of explicit
   elementary Mellin transforms for two independent test functions.\\[4pt]
Moment generators
 & 0 & 3
 & Anchored quadrature, explicit polylogarithmic primitives, and a Taylor
   sum using 65 exact moment coefficients in each comparison.\\[4pt]
\midrule
Total & 1,050 & 72 & Finite algebra and numerical diagnostics have
                         different evidentiary roles.\\
\bottomrule
\end{tabular}
\end{center}

The 590 kernel checks use rational arithmetic only.  Positive
fourth-power exponential nodes and quarter-integral endpoints permit
exact tests even when $a-b$ is nonintegral.  The 460 moment checks are
finite identities over rational symbolic coefficient rings.  In
particular, the finite summation-by-parts tests keep the two tail
constants symbolic.

The transverse diagnostics use 65 decimal working digits and include
the discrepancy \eqref{audit:continued}--\eqref{audit:literal} as a
deliberate boundary control.  A separate 80-digit test of
\eqref{tr:Mjet} uses two quadratic-polynomial exponential functions,
one with nonreal coefficients and decay parameter.  For each function,
the centers $m=0,1,3$ and derivative orders $r=1,2,3,4$ compare the
subtracted logarithmic integral with a numerical derivative of its
known elementary Mellin transform.  All 24 cases satisfy the scripted
$10^{-50}$ tolerance; the largest observed discrepancy is about
$3.10\times10^{-79}$.  These functions test the normalization lemma
independently of the relative harmonic kernel.

The moment diagnostics use 95 working
digits, with orders through $(6,8)$ and moment indices through $M=7$.
They evaluate a direct finite Hurwitz-zeta head and a separately
assembled centered asymptotic tail at $(N,K)=(48,32)$ and $(64,36)$.
The largest observed discrepancy at the larger cutoff is about
$1.16\times10^{-63}$.  The three generator comparisons include a
nonreal parameter; their largest recorded Taylor-sum discrepancy is
about $1.08\times10^{-89}$.  The 65 Taylor coefficients are exact
inputs to those three numerical comparisons, not 65 additional
independent exact tests.

These residuals are observed floating-point discrepancies, not rigorous
error bounds or interval enclosures.  Working precision alone is not
an accuracy guarantee.  The recorded diagnostics cover the stated
examples; they do not constitute numerical quadrature of every
multidimensional integral in Theorem~\ref{tr:integral}.  Uniform
integrability, simultaneous branch cancellation, and continuation at
arbitrary orders are established by the proofs, not by a finite test
count.  The delivery includes neither a proof-assistant formalization
nor an external referee report.

\subsection{Status of the broader conjectures}

The contribution ledger records R7 as answered in the specified
arbitrary-depth zero-decoration family, R8 as answered for one marked
index and all of its nonpositive-center derivatives, and R1 as answered
for the full two-factor even-tail sector.  Extra harmonic factors,
several independently marked indices, and additional global spectral
derivative reductions remain separate questions.

The Gaussian $S_6$ and revised $S_8$ candidates retain their
conjectural status \cite{Repository,IncomingParity}.  No new functional
identity or exact period certificate for either candidate is supplied
here.  Existing nonmembership results concern their explicitly named
formal relation spaces and must not be read as period nonvanishing
theorems.  The next research questions in
Section~\ref{research:section} begin at these stated boundaries.
